\documentclass[11pt]{article}

\usepackage[margin=1in]{geometry}
\usepackage[T1]{fontenc}
\usepackage{lmodern}
\usepackage{microtype}
\usepackage{amsmath,amssymb,amsthm,mathtools}
\usepackage[dvipsnames]{xcolor}
\usepackage[most]{tcolorbox}
\usepackage{tikz}
\usepackage{enumitem}
\usepackage{tocloft}
\usepackage[hidelinks]{hyperref}

\definecolor{courseblue}{HTML}{1F4E79}
\definecolor{lightblue}{HTML}{EAF2F8}
\definecolor{lightgreen}{HTML}{EAF7EA}
\definecolor{lightpink}{HTML}{FCE8EF}
\definecolor{lightyellow}{HTML}{FFF8D6}
\definecolor{warningyellow}{HTML}{D6A000}
\definecolor{exercisebg}{HTML}{D6E6F5}
\definecolor{exerciseframe}{HTML}{1F4E79}
\definecolor{examplebg}{HTML}{F4FAFE}
\definecolor{exampleframe}{HTML}{85C1E9}
\definecolor{darkgray}{HTML}{333333}

\setlength{\parindent}{0pt}
\setlength{\parskip}{0.65em}
\setlist[itemize]{topsep=0.25em,itemsep=0.15em}
\renewcommand{\cftsecleader}{\cftdotfill{\cftdotsep}}
\renewcommand{\cftsubsecleader}{\cftdotfill{\cftdotsep}}

\theoremstyle{plain}
\newtheorem{theorem}{Theorem}[section]
\newtheorem{proposition}[theorem]{Proposition}
\newtheorem{lemma}[theorem]{Lemma}
\newtheorem{corollary}[theorem]{Corollary}
\theoremstyle{definition}
\newtheorem{definition}[theorem]{Definition}
\newtheorem{example}[theorem]{Example}
\newtheorem{exercise}[theorem]{Exercise}
\theoremstyle{remark}
\newtheorem*{remark}{Remark}
\newtheorem*{fact}{Fact}
\newtheorem*{warning}{Warning}

\tcolorboxenvironment{definition}{
  enhanced,
  colback=lightgreen,
  colframe=ForestGreen,
  boxrule=0.7pt,
  arc=2pt,
  left=6pt,
  right=6pt,
  top=5pt,
  bottom=5pt
}

\tcolorboxenvironment{warning}{
  enhanced,
  colback=lightyellow,
  colframe=warningyellow,
  boxrule=0.7pt,
  arc=2pt,
  left=6pt,
  right=6pt,
  top=5pt,
  bottom=5pt
}

\tcolorboxenvironment{exercise}{
  enhanced,
  colback=exercisebg,
  colframe=exerciseframe,
  boxrule=1pt,
  arc=2pt,
  left=6pt,
  right=6pt,
  top=5pt,
  bottom=5pt
}

\tcolorboxenvironment{example}{
  enhanced,
  colback=examplebg,
  colframe=exampleframe,
  boxrule=0.7pt,
  arc=2pt,
  left=6pt,
  right=6pt,
  top=5pt,
  bottom=5pt
}

\numberwithin{equation}{section}

\newcommand{\be}{\begin{eqnarray}}
\newcommand{\ee}{\end{eqnarray}}
\newcommand{\C}{\mathbb{C}}
\newcommand{\R}{\mathbb{R}}
\newcommand{\N}{\mathbb{N}}

\DeclareMathOperator{\Hol}{\mathcal{H}\mathrm{ol}}
\DeclareMathOperator{\Real}{Re}
\DeclareMathOperator{\Imag}{Im}
\newcommand{\A}{\mathcal{A}}
\newcommand{\CT}{\C[[T]]}
\newcommand{\CLT}{\C((T))}
\tcbset{breakable}
\begin{document}

\begin{center}
  \colorbox{lightblue}{%
    \begin{minipage}{0.94\textwidth}
      \vspace{0.55em}
      \begin{tabular*}{\textwidth}{@{\extracolsep{\fill}}lr}
        {\color{courseblue}\bfseries\LARGE Lecture 8} &
        {\color{courseblue}\bfseries\LARGE MATH 411}\\[0.35em]
        \textbf{Fall 2026} & \textbf{Date: Sep 21, 2026}\\
        & \textbf{Instructor: Manish Patnaik}
      \end{tabular*}
      \vspace{0.45em}
    \end{minipage}}
\end{center}

\setcounter{tocdepth}{1}
\tableofcontents
\vspace{0.45em}

% Handwritten page 1.
\section{What does it mean to plug a complex number into a formal power series?}

Let $f\in\CT$ be a formal power series,
\be
 f=\sum_{n\ge0}a_nT^n,\qquad a_n\in\C.
\ee
The symbol $T$ is an indeterminate: so far nothing has been substituted for
it. The question driving this lecture is: \emph{what sense does it make to
``plug in'' a number $z\in\C$?}

The answer goes through series of functions. Let $S$ be a set and let
$f_n\colon S\to\C$ be functions on $S$; we consider
\be
 \sum_{n} f_n.
\ee
The case we actually care about is
\be
 f_n(z)=a_nz^n,\qquad S=D(0,r),
\ee
so that the partial sums are the truncations of the power series.

\begin{definition}[The three modes of convergence]
Write $S_k=\sum_{n=0}^{k}f_n$ for the $k$-th partial sum. We say that
$\sum f_n$
\begin{itemize}
\item converges \emph{pointwise} if, for each $z\in S$, the sequence of
complex numbers $\{S_k(z)\}$ converges;
\item converges \emph{uniformly} if $\{S_k\}$ converges with respect to the
sup norm $\|\cdot\|_S$;
\item converges \emph{absolutely} if, setting
$\displaystyle V_k(z)=\sum_{n=0}^{k}|f_n(z)|$, the sequence $\{V_k(z)\}$
converges for each $z\in S$.

\end{itemize}
\end{definition}

% Handwritten page 2.
What we care about is the implication
\be
 \sum |a_n|r^n<\infty
 \quad\Longrightarrow\quad
 \sum a_nz^n\text{ converges uniformly and absolutely on }
 \overline{D(0,r)},
\ee
which is an instance of the comparison test that we proved last time:


\begin{theorem}[Comparison test]
\label{thm:comparison}
Let $c_n\ge0$ be real numbers with $\sum c_n$ convergent, let $S\subseteq\C$,
and let $\{f_n\}$ be functions on $S$ with
\be
 \|f_n\|_S\le c_n\qquad\text{for all }n.
\ee
Then $\sum f_n$ converges uniformly and absolutely on $S$.
\end{theorem}
%
%\begin{proof}
%For $m>k$,
%\be
% \left\|\sum_{n=k+1}^{m}f_n\right\|_S
% \leq \sum_{n=k+1}^{m}\|f_n\|_S
% \leq \sum_{n=k+1}^{m}c_n.
%\ee
%Since $\sum c_n$ converges, its tails tend to $0$. Thus the partial sums of
%$\sum f_n$ are uniformly Cauchy. At each $z\in S$ they therefore converge to
%some value $f(z)$, and the same tail estimate shows that the convergence to
%$f$ is uniform. Moreover, for each $z\in S$,
%\be
% \sum_n|f_n(z)|\leq\sum_n c_n<\infty,
%\ee
%so the convergence is absolute at every point of $S$.
%\end{proof}

Indeed, the application we have in mind is the following:

\begin{corollary}[Application]
\label{cor:closed-disc}
Suppose $\sum|a_n|r^n$ converges for some $r>0$. Then $\sum a_nz^n$
converges uniformly and absolutely on the \emph{closed} disc
$\overline{D(0,r)}$.
\end{corollary}

\begin{proof}
Let $S=\overline{D(0,r)}=\{z\in\C:|z|\le r\}$ and put $f_n(z)=a_nz^n$. Then
\be
 \|f_n\|_S\le|a_n|\,r^n,
\ee
so Theorem~\ref{thm:comparison} applies with $c_n=|a_n|r^n$.
\end{proof}

% Handwritten page 3.
\section{The radius of convergence}

For a power series $\sum a_nz^n$, let
\be
 T=\{0\}\cup
 \{s>0 \;:\; \textstyle\sum|a_n|s^n\text{ converges}\},
\ee
and let $r_0=\sup T\in[0,\infty]$.

\begin{definition}[Radius of convergence]
The extended nonnegative number $r_0$ above is called the \emph{radius of
convergence} (R.O.C.) of $\sum a_nz^n$. Since $0\in T$, every power series
has a well-defined radius of convergence. The extreme cases are:
\begin{itemize}
\item $r_0=\infty$: the series converges absolutely at every $z\in\C$ and
uniformly on every closed disc;
\item $r_0=0$: the series does not converge anywhere except at $z=0$.
\end{itemize}
\end{definition}

\begin{proposition}
The series $\sum a_nz^n$ converges absolutely at every point of $D(0,r_0)$
and uniformly on every closed disc $\overline{D(0,s)}$ with $s<r_0$. It
diverges for $|z|>r_0$.
\end{proposition}

\begin{proof}
If $0\leq s<r_0$, the definition of the supremum gives some $q\in T$ with
$s<q$. For $|z|\leq s$,
\be
 |a_nz^n|\leq |a_n|s^n\leq |a_n|q^n,
\ee
so the comparison test gives uniform and absolute convergence on
$\overline{D(0,s)}$.

Now suppose $|z|>r_0$ and $\sum a_nz^n$ converges. Its terms are bounded, so
$|a_nz^n|\leq M$ for some $M>0$. Choose $q$ with $r_0<q<|z|$. Then
\be
 \sum_{n\geq0}|a_n|q^n
 \leq M\sum_{n\geq0}\left(\frac q{|z|}\right)^n<\infty,
\ee
which says $q\in T$, contradicting $q>\sup T=r_0$.
\end{proof}

\begin{warning}
When $0<r_0<\infty$, nothing above says what happens on the circle
$|z|=r_0$ itself. The behaviour there genuinely depends on the series: it may
converge at every boundary point, at none, or at some but not others.
\end{warning}

% Handwritten page 4.
\section{Interlude: \texorpdfstring{$\limsup$}{limsup}}

Hadamard's theorem, stated below, computes $r_0$ directly from the
coefficients, but it is phrased using $\limsup$. So: what does $\limsup$
mean?

\begin{definition}[Accumulation point]
Let $(t_n)_{n\geq0}$ be an infinite sequence of nonnegative real numbers. An
\emph{accumulation point} of $(t_n)_{n\geq0}$ is a point $t\in\R$ such that
every interval around $t$ contains $t_n$ for infinitely many indices $n$.
\end{definition}

\begin{fact}
Every bounded infinite sequence in $\R$ has a convergent subsequence and hence
an accumulation point.
\end{fact}

\begin{definition}[$\limsup$]
Suppose the infinite sequence $(t_n)_{n\geq0}$ is bounded, and let $S$ be the
set of all its accumulation points. Then
\be
 \limsup_{n} t_n=\text{l.u.b.}(S).
\ee
\end{definition}

A practical way of viewing $\limsup$ is that it is the number
$\lambda\in\R$ characterised by: for every $\epsilon>0$,
\begin{itemize}
\item[(A)] there exist only \emph{finitely many} $n$ with
$t_n\ge\lambda+\epsilon$;
\item[(B)] there exist \emph{infinitely many} $n$ with
$t_n\ge\lambda-\epsilon$.
\end{itemize}

% Handwritten page 5.
\begin{definition}[Convention]
If the infinite sequence $(t_n)_{n\geq0}$ is unbounded, we set
$\limsup_n t_n=\infty$.
\end{definition}

\begin{exercise}
Let $(t_n)_{n\geq0}$ and $(s_n)_{n\geq0}$ be bounded infinite sequences of
nonnegative real numbers and set $t=\limsup_n t_n$, $s=\limsup_n s_n$. Verify:
\begin{itemize}
\item $\displaystyle\limsup_n(t_n+s_n)\le t+s$;
\item $\displaystyle\limsup_n(t_ns_n)\le t\cdot s$;
\item if $\displaystyle\lim_{n\to\infty}t_n$ exists, then
$\displaystyle\limsup_n t_n=\lim_{n\to\infty}t_n$.
\end{itemize}
\end{exercise}

\section{Hadamard's criterion}

\begin{theorem}[Hadamard]
\label{thm:hadamard}
Let $r_0$ be the radius of convergence of $\sum a_nz^n$. Then
\be
 \frac{1}{r_0}=\limsup_{n}|a_n|^{1/n}.
\ee
Here we use the conventions $1/0=\infty$ and $1/\infty=0$.
\end{theorem}

\begin{proof}
Let $t=\limsup_{n\to\infty}|a_n|^{1/n}$ and assume $0<t<\infty$ (the extreme
cases are treated in the remark below). We must show that $\sum a_nz^n$
converges for $|z|<1/t$ and diverges for $|z|>1/t$; this identifies the
radius of convergence as $r_0=1/t$.

% Handwritten page 6.
\textbf{Convergence.} By characterisation (A) above, for any $\epsilon>0$
there are only finitely many $n$ with
\be
 |a_n|^{1/n}\ge t+\epsilon,
 \qquad\text{i.e.}\qquad |a_n|\ge(t+\epsilon)^n.
\ee
Equivalently,
\be
\label{eq:club}
 |a_n|<(t+\epsilon)^n\qquad\text{for almost all }n.
\ee
Let $z\in\C$ with $|z|<1/t$. Choose $\epsilon>0$ small enough that
\be
 |z|(t+\epsilon)<1,\qquad\text{i.e.}\qquad |z|<\frac1{t+\epsilon}.
\ee
Then, using~\eqref{eq:club} for all but finitely many terms (the finitely
many exceptions contribute a constant $c$),
\be
 \sum|a_n||z|^n\le c+\sum\bigl((t+\epsilon)|z|\bigr)^n,
\ee
and the latter is a convergent geometric series since
$(t+\epsilon)|z|<1$. By comparison, $\sum|a_n||z|^n$ converges.

\textbf{Divergence.} Assume $|z|>1/t$ and choose $\epsilon$ with
$0<\epsilon<t$ such that $|z|=\dfrac1{t-\epsilon}$.
% Handwritten page 7.
For this $\epsilon>0$, characterisation (B) gives infinitely many $n$ with
\be
 |a_n|^{1/n}\ge t-\epsilon,\qquad\text{hence}\qquad |a_n|\ge(t-\epsilon)^n.
\ee
Therefore, for infinitely many $n$,
\be
 |a_nz^n|\geq (t-\epsilon)^n\cdot\frac1{(t-\epsilon)^n}=1.
\ee
In particular, the terms $a_nz^n$ do not tend to $0$, so $\sum a_nz^n$
diverges.
\end{proof}

\begin{remark}
If the limsup in Hadamard's criterion is zero, then the series converges
everywhere, so $r=\infty$. If the limsup is infinite, then $r=0$.
\end{remark}

\begin{example}
Consider
\be
 \sum_{n\ge0}\frac{z^n}{n!},\qquad a_n=\frac1{n!}.
\ee
To apply Hadamard, look at $\displaystyle\limsup_n\left(\frac1{n!}\right)^{1/n}$.
By Stirling's formula, $n!=n^ne^{-n}u_n$ with
$\displaystyle\lim_{n\to\infty}u_n^{1/n}=1$. Plugging this in gives
\be
 \limsup_{n}\left(\frac1{n!}\right)^{1/n}=0,
\ee
so the radius of convergence is $\infty$: the exponential series converges
on all of $\C$, as we saw last lecture.
\end{example}

% Handwritten page 8.
\section*{Next time}
Let $f(z)=\sum a_nz^n$ have radius of convergence $r$. What is the radius of
convergence of the formally differentiated series
\be
 \sum_{n\geq1} n\,a_nz^{n-1}\,?
\ee
It turns out to be $r$ again --- and this is the first step towards showing
that a convergent power series defines a $\C$-differentiable function.

\end{document}
